% Bruchrechnung · Addieren/Subtrahieren · Prüfungs-Fokus MSA – bank/bruchrechnung/e1.jsonl (zusammenbau v0.6, Prüfungs-Fokus)
\einheitenkopf[e1]{Einheit 1 · Brüche addieren und subtrahieren}
\zweigzeile{ab Kl. 5, je nach Buch bis Kl. 6 (am Gymnasium ab Kl. 6) · P10 ×10}
\setcounter{aufgabe}{0}

% Anlauf (ohne Prüfkennung)
\begin{aufgabe}{Addieren/Subtrahieren – Anlauf}
\begin{teile}
\teil Färbe $\frac{1}{6}$, dann weitere $\frac{4}{6}$ – wie viel ist gefärbt?

\bruchrechteck{0}{6}
\leerfeld
\teil $\frac{2}{11} + \frac{5}{11}$
\teil $\frac{5}{14} + \frac{3}{14}$
\end{teile}
\end{aufgabe}

\begin{aufgabe}{Addieren/Subtrahieren – Prüfungsaufgaben}
\begin{teile}
\steil Ein Glücksrad hat 10 gleich große Felder: 5 rot, 3 blau, 2 gelb. Es wird zweimal gedreht. Die Pfade „zweimal rot“, „zweimal blau“ und „zweimal gelb“ haben die Wahrscheinlichkeiten $\frac{1}{4}$, $\frac{9}{100}$ und $\frac{4}{100}$. Wie groß ist die Wahrscheinlichkeit für zweimal dieselbe Farbe? Gib sie gekürzt an. \hfill (P14*)
\steil In einer Tüte sind 10 Bonbons: 6 Kirsche, 3 Zitrone, 1 Apfel. Du ziehst eins, legst es zurück und ziehst noch eins. Die Pfade „zweimal Kirsche“, „zweimal Zitrone“ und „zweimal Apfel“ haben die Wahrscheinlichkeiten $\frac{9}{25}$, $\frac{9}{100}$ und $\frac{1}{100}$. Wie groß ist die Wahrscheinlichkeit für zweimal dieselbe Sorte? Gib sie gekürzt an. \hfill (P14*)
\steil Ein Los bestimmt die Reihenfolge von 9 Kindern beim Staffellauf. Die Wahrscheinlichkeiten, dass Mia als Erste, als Zweite oder als Dritte gezogen wird, sind $\frac{1}{9}$, $\frac{8}{72}$ und $\frac{56}{504}$. Wie groß ist die Wahrscheinlichkeit, dass Mia unter den ersten drei ist? Gib sie gekürzt an. \hfill (P15*)
\steil Beim Vorlesewettbewerb lost die Lehrerin die Reihenfolge von 12 Kindern aus. Die Wahrscheinlichkeiten, dass Jan als Erster, als Zweiter oder als Dritter liest, sind $\frac{1}{12}$, $\frac{11}{132}$ und $\frac{110}{1320}$. Wie groß ist die Wahrscheinlichkeit, dass Jan unter den ersten drei ist? Gib sie gekürzt an. \hfill (P15*)
\steil Eine Münze wird dreimal geworfen. Die Wahrscheinlichkeit für genau zweimal Zahl ist $\frac{3}{8}$, für dreimal Zahl $\frac{1}{8}$. Wie groß ist die Wahrscheinlichkeit für mindestens zweimal Zahl? Gib sie gekürzt an. \hfill (P20*)
\steil Sara wirft drei Freiwürfe. Die Wahrscheinlichkeit für genau zwei Treffer ist $\frac{27}{64}$, für drei Treffer ebenfalls $\frac{27}{64}$. Wie groß ist die Wahrscheinlichkeit für mindestens zwei Treffer? Gib sie gekürzt an. \hfill (P20*)
\steil Tim hat vier mögliche Codes für sein Fahrradschloss und probiert sie nacheinander, ohne einen zu wiederholen. Die Wahrscheinlichkeit, dass der erste Versuch passt, ist $\frac{1}{4}$; dass erst der zweite passt, ist $\frac{3}{12}$. Wie groß ist die Wahrscheinlichkeit, dass das Schloss spätestens beim zweiten Versuch aufgeht? Gib sie gekürzt an. \hfill (P17*)
\steil Nina weiß, dass ihr Passwort eines von acht möglichen ist, und probiert sie nacheinander, ohne eines zu wiederholen. Die Wahrscheinlichkeit, dass der erste Versuch passt, ist $\frac{1}{8}$; dass erst der zweite passt, ist $\frac{7}{56}$. Wie groß ist die Wahrscheinlichkeit, dass sie spätestens beim zweiten Versuch drin ist? Gib sie gekürzt an. \hfill (P17*)
\teil Zwei Kreisel werden gedreht. Die Wahrscheinlichkeit, dass beide Rot zeigen, ist $\frac{4}{30}$. Wie groß ist die Wahrscheinlichkeit, dass nicht beide Rot zeigen? Gib sie gekürzt an. \hfill (P26F)
\teil Ali zieht an zwei Losbuden je ein Los. Die Wahrscheinlichkeit für zwei Gewinne ist $\frac{6}{48}$. Wie groß ist die Wahrscheinlichkeit, dass er nicht zweimal gewinnt? Gib sie gekürzt an. \hfill (P26F)
\teil Ein Glücksrad hat 10 gleich große Felder mit den Zahlen 1 bis 10. Die Wahrscheinlichkeit für 8, 9 oder 10 ist $\frac{3}{10}$. Wie groß ist die Wahrscheinlichkeit für eine Zahl von 1 bis 7? \hfill (P14)
\teil Ein Spielwürfel hat 12 Flächen mit den Zahlen 1 bis 12. Die Wahrscheinlichkeit für 1 oder 12 ist $\frac{2}{12}$. Wie groß ist die Wahrscheinlichkeit für weder 1 noch 12? Gib sie gekürzt an. \hfill (P14)
\end{teile}
\end{aufgabe}

\medskip
\uebersichtskasten{Gleicher Nenner: Zähler plus (oder minus), Nenner bleibt. \quad 3/7 + 2/7 = 5/7 \\ Verschiedene Nenner: erst gleichnamig machen (erweitern), dann Zähler rechnen. \quad 1/2 + 1/6 = 3/6 + 1/6 = 4/6 = 2/3 \\ Ergebnis kürzen; über ein Ganzes: 9/5 = 1 4/5.}
